\section{Why triple relations alone do not give small separators}
\label{sec:coordinate-separator-obstruction}

The balanced compatibility-separator hypothesis of
Section~\ref{sec:euler-coordinate-factors} needs more than the local
three-column relation, rank $t-1$ and polynomial encoding.  The following
explicit family shows this for abstract rational cones.  It is not a family
of native triangulations and has no prescribed canonical Euler function.
Consequently it proves no lower bound for unknot recognition or its resolving
parameter.

\begin{proposition}[An explicit family with no small coordinate separator]
For every $t\ge3$ there is an integer $(t-1)\times3t$ matrix $C$, partitioned
into $t$ triples summing to zero, with $O(t\log t)$-bit entries, such that:
\begin{enumerate}
\item Every set of at most $t-1$ columns using at most two from each triple
is independent.  In particular $\operatorname{rank}C=t-1$.
\item After restricting any set of $s$ triples to arbitrary singletons,
the retained column matroid has one component of size $3t-2s$.
\item If $s\le\lfloor(t+1)/2\rfloor$, the restricted nonnegative kernel
contains a vector strictly positive on every retained coordinate.  Its
nonnegative cone has no nontrivial literal coordinate product, even after
exact implicit-zero deletion.
\end{enumerate}
For balance factor $2/3$, the smallest singleton separator has size
$\lceil t/2\rceil$, with or without implicit-zero deletion.
\end{proposition}

\begin{proof}
Put $r=t-1$ and define
\[
 v(a)=(1,a,\ldots,a^{r-1})^{\mathsf T},\qquad
 u_i=v(2i-1),\quad v_i=v(2i),\quad w_i=-u_i-v_i
 \quad(1\le i\le t).
\]
Use these triples as the columns of $C$, in increasing group order.
For any $r$ chosen columns using at most two per group, expand its determinant
multilinearly.  Within a two-column group, the three possible pairs reduce
with fixed signs:
\[
 (u_i,v_i)\mapsto +(u_i,v_i),\qquad
 (u_i,w_i)\mapsto -(u_i,v_i),\qquad
 (v_i,w_i)\mapsto +(u_i,v_i).
\]
Here the notation describes their contribution to the determinant with all
other columns fixed.  A singleton $w_i$ contributes the negative sum of the
choices $u_i$ and $v_i$.  Every resulting term therefore has the same overall
sign.  Its moment-curve nodes are distinct and increasing: the two nodes
of a group are adjacent and precede every node of the next group.  Each
remaining determinant is the strictly positive Vandermonde product
\[
                      \prod_{a<b}(x_b-x_a).
\]
There is at least one term and none can cancel, so the original minor is
nonzero.  A smaller admissible set extends to such an $r$-element set, proving
the first assertion.

Every transversal with one column per group has $t=r+1$ columns and is
dependent, while every proper subset is independent.  It is therefore a
circuit.  In any singleton-restricted mask, choose one retained column in
each group to form a transversal $T_0$.  All its columns lie in one matroid
component.  Any other retained column replaces its group's member of $T_0$
to give another circuit sharing $t-1>0$ columns with $T_0$.  Thus every
retained column is in that same component.  The rank is still $r$ and the
nullity is $2t+1-2s$.

For the third assertion, put value one in every fixed singleton coordinate.
There are $g=t-s$ untouched full groups.  If
$s\le\lfloor(t+1)/2\rfloor$, then $2g\ge r$, so we can choose $r$ columns
from these groups using at most two per group.  Their matrix is invertible.
Solve for signed coefficients on those columns that cancel the sum of the
fixed singleton columns, setting the other untouched coordinates to zero.
This produces a signed kernel vector $z$ whose singleton values are one.
In each untouched group, add
\[
              \lambda_i(1,1,1),\qquad
              \lambda_i=1+\max(0,-z_j:j\text{ in group }i).
\]
The group sum relation keeps the vector in the kernel and makes every
untouched coordinate strictly positive.  Singleton coordinates remain one.
Thus there are no implicit zeros.  The cone spans the signed kernel by the
full-support proposition of Section~\ref{sec:euler-coordinate-factors};
the connected matroid excludes every nontrivial coordinate cone product.

The single remaining factor has $3t-2s$ columns.  To reduce it to at most
$(2/3)3t=2t$ requires $s\ge\lceil t/2\rceil$.  Below this threshold full
support prevents implicit-zero deletion from helping.  At the threshold,
any choice of groups suffices just by removing the competing columns, and
full support still holds.  This proves the exact separator size.
\end{proof}

The construction is explicit and polynomial in its output length.  Each
moment entry has at most $(t-2)\log_2(2t)+O(1)$ bits, and adding the two
entries for $w_i$ adds at most one bit.  The full matrix therefore has
$O(t^3\log t)$ bits.  Cramer's rule applied to the signed solve in the proof
also gives a full-support rational vector of polynomial encoding length;
the right-hand side is a sum of at most $t$ input columns.  The subsequent
group shifts preserve that polynomial bound.  No genericity assumption or
search over exponentially many determinants is needed to construct the
matrix.

More generally, without implicit-zero deletion, any fixed balance factor
$0<\alpha<1$ requires
\[
                s\ge\left\lceil\tfrac32(1-\alpha)t\right\rceil.
\]
If this number exceeds $t$, the specified singleton/factor mechanism cannot
achieve that balance.  With implicit-zero deletion allowed, a necessary
bound is
\[
 s\ge\min\left\{\lfloor(t+1)/2\rfloor+1,
                 \left\lceil\tfrac32(1-\alpha)t\right\rceil\right\},
\]
which is still $\Omega(t)$ for fixed $\alpha<1$.  Outside the proved
full-support range, signed connectivity alone cannot justify the stronger
cone conclusion.  Indeed fixing all groups to their $u_i$ columns makes the
nonnegative kernel zero, since every entry in the first row is positive,
even though those columns still form a signed circuit.

This family satisfies the tempting algebraic premises uniformly, including
full support after every sufficiently small singleton restriction.  Thus
those premises cannot establish polylogarithmic balanced coordinate
separators.  A theorem for the native sources must use further restrictions
on matrices realized by authenticated knot-exterior triangulations, or use
their canonical Euler functions to close branches without an algebraic
split.  The family supplies neither such a realization nor an Euler-positive
ray.  It also says nothing about non-coordinate decompositions or the cost
of directly deciding these abstract cones; a connected matching cone can
still be rejected by an easy Euler query.
